\documentclass{jsarticle}
\usepackage[dvipdfmx,final]{graphicx}
\usepackage{amsmath}
\usepackage{amssymb}
\begin{document}
\title{行動による人が持っている情報ネットワークの \\
可能性の発見と再構築、情報空間 \\
によるモデル}
\author{苫米地博士の論文}
\date{}
\maketitle

      初速度のドーパミンの量からの行動の補助から、終わりの精神のドーパミン
  の量の調節、相加相乗平均から、始まりと終わりの精神のドーパミンの
  調節による人の行動の領域による予知
  $$ x \log x = \left(v_0 + {1 \over at_n}\right)V, V_x = e^f + {1 \over e^f} $$
     経路積分によるエントロピー不変量から
  このドーパミンの量から人の行動傾向が決まる。
  基礎代謝の持続できる距離と精神の介入と補助、言語の発生と抑制による経路
  パターンの組み合わせを再構築する。
  $$ ||dx^2|| = \int {\mathrm{exp}\left({1 \over \sqrt{2 \tau q}}L(x)
  \right) + \mathrm{O}(N^{-1})}d \tau   $$
  $$ p = {v \over 2m}, \hbar = {p \over 2mv}, \lambda = h \nu $$
     これらの式から、言語の神経ネットワーク
  による単体クラスの組み合わせパターンとその対象の人の領域による
  人が持っている情報空間の傾向の可能性のポテンシャルエネルギーの
  予知を人工知能の正規表現で抽出する。
   
  $$ {d \over d \gamma}\Gamma = \Gamma^{{\gamma}^{'}} = e^{-f}  $$
  $$ C = \int \left(\int {1 \over x^s}dx + \log x \right)\mathrm{dvol} $$
     オイラーの定数は、ゼータ関数とゼータ関数自体のエントロピーを
  体積積分した結果でもある。このオイラーの定数は、ガンマ関数の
  大域的微分方程式からも導出できる。このガンマ関数は
  大脳基底核のエントロピー不変量の強度も表している。
  このガンマ関数の大域的微分多様体による、ブレインインターフェース
  での人体のデータ計測もこのJones多項式とラプラス方程式からの
  ３次元多様体のエントロピー不変式からも、言語とイメージの
  大脳基底核からのデータ抽出も可能と言える。
  これらのデータは、口腔の周りの熱エネルギーの流れと動きで、音階と音層、発声器との
  合成でわかる。速読していてもわかる。目の瞳孔と頭頂葉による熱流体でもわかる。
  ブローカー中枢の発生器としての機能と、ウェルニッケ中枢の音を感知する役割を
  利用して、外部に観測器として、温度計とカメラで解決している。
  脳の内部に埋め込むと、統合失調症や、躁うつ病の人に、思考伝播と考察伝播が出現して、
  妄想状態になるので、外部観測で実用化を解決したらしい。


　　ブレインインターフェースとその応用、機知の仕組み



　ブレインインターフェースは、機械において、音声を文字列に表示するには、
まず、音声を音階の7階層に分解して、その上に、大脳基底核を経由する熱エネルギーの共鳴エネルギーを使って、
唇の動きのデータを外部で電気信号に変えて、この2種類のデータから、音声を文字列に表示する。
頭頂葉から大脳基底核を経由する体感触覚による熱感気のデータを使って、そのデータをイメージと言語と文字列に変えて、画像処理を使って、動的に映像やページを処理する。例えば、Wordのページをめくったりする。
ASを頭頂葉から大脳基底核を経由する熱エネルギーのデータを電気信号に変えて、
機知の仕組みを担う第32，33野を経由するドーパミンがそれぞれの役割をする共振回路の代わりをする
熱エネルギーの共鳴を使って、血液のリボ核酸のヘリコバクターが共鳴する共振回路の役割の代わりの共鳴エネルギー
に外部でのカメラと温度計を介して接続して、ASを動かす。

$$ H_n=\sigma(E_n )×\sigma(K_n), e_n=f^{-1}(x)xf(x)$$
$$e_n=\mathrm{ker}f \oplus \mathrm{im}f,H_n^2=|(|ds^2 |)|,=H_n×K_n $$

$$E_n=σ(K_n )×σ(H_n ),π(σ_1,σ_2 )=  {d \over df} F(x) $$
$$ x\Gamma(x) =\pi(\sigma_1, \sigma_2 )-  (8 \pi G({p \over c^3} + 
{V \over S})) $$
$$=2\int  ||\sin2x||^2 d \tau $$
$$ =4 \int ||\sin x  \cos x||^2 d \tau \to 4V = g_{ij}^2 $$

   固有エントロピー = 原子の固有エントロピー - ３次元多様体のエントロピー


   ブレインインターフェースの大脳基底核の固有エントロピー値





     Jones多項式は、この熱エネルギーのエントロピー不変量による、音階の音層が、
  周期関数になっていることと、この周期関数がエネルギー流体として、音が、血流断層を抵抗値としての
  強度を測れて、唇がどの音の言葉の発声に近いかを、この2種類で合成して、音声認識が出来ている。
  温度計とカメラで解決されたらしい。脳に埋め込まれるのは、ずっと前の人で、今の人は大丈夫です。
  アップルとマイクロソフトとグーグルと富士通、NECに助けられました。
  荒木先生が、特効薬の心理薬をくださり、完治しました。
  蔵王病院は、そのような心理薬をたくさんのデータとして、所有しております。
  私みたいな、困った症状の人は、荒木先生か蔵王病院、広島大学医学部、富山大学医学部に行かれたら、
  最高でも12年で治ります。彩さんの論文と、苫米地博士の論文をアカシックレコードと瞑想法で見られたら、
  もっと治ります。
  この2人は、ノーベル経済学賞とノーベル物理学賞とノーベル医学・生理学賞抜擢です。
   統合失調症、躁うつ病が治ります。論文見るべきです。条件は、アカシックレコードと瞑想法とです。
   勝間和代先生も合わせると大丈夫です。蔵王病院の石川博先生と河相和昭先生、脇本和博先生、葉子先生、
   奥田伴枝先生、光の丘病院の馬屋原先生、小林心理士先生、石岡芳隆先生、広島大学医学部の横田則夫先生、
   梶山先生、辻先生、後藤先生、萬谷智之先生、中村先生、渡辺千種先生、小鶴先生、
   仁風荘の福場先生、鮎沢香織先生、
   富山大学医学部の田尻先生、
   脇本先生、田辺先生が、私は推薦します。
   その後に、デイケアささっといってほしいらしく、仕方ないです。

      SRS研究所の田中先生、栗田先生、フォトリーディングのポール・シーリィ先生、園先生、勝間和代先生、
   瞬読の山中先生、うちのユウちゃんも推します。
   GSR速読法もいいです。スポーツ速読法もいいです。王様の速読の速読耳もいいです。
   川島先生もいいです。
   どれか一つの速読法にした方がいいらしく、そうでないと、音読とフィフティーフィフティーになり、後ろ指指されます。

\title{Jones多項式からの各株式発行からの需要と供給、\\
株価の分配関数、多項式からナスダックへの金融機構の \\
株式発行のマクロとミクロの波及と流れ}
\author{高島彩さんの成蹊大学での経済理論の研究論文}

   Jones多項式から各会社の株価への値打ちから日経平均株価への
   株価の価値の波及、東証株価指数Topixの全銘柄からの全体の会社の
   指数からの平均株価の決定量、そして、
   アメリカの株式市場でのダウ平均株価の終値の収束
   (戦略、reco level theory,株式会社への)全体（株式会社
   への）需要と供給。始値から終値への分配、このJones多項式の出力から
   終値への各同値類と剰余類からの部分群論の、各平均株価の推移率の予測、
   この応用は、１９２９年の世界恐慌がニューヨークダウ平均株価の暴落が、
   このJones多項式からの流れが、終値で破綻して、金融の共鳴力が、
   干渉してしまい、カタストロフィー現象として、金融の流れが乱れた結果で
   あると、この論文から分析できる。
   $$ \Hat{V} = {1 \over \sqrt{t}} \vec{V} - {1 \over t}\Hat{V} $$
   $$ \Hat{V} = \sqrt{t}\Hat{V} - t \Hat{V} $$
   $$ \Hat{V}_k = \sum_{\vec{\sigma}} \langle K| \vec{\sigma}\rangle
   \left(\sqrt{t} + {1 \over \sqrt{t}} \right)^{||\vec{\sigma}||} $$
   $$ {1 \over \sqrt{t}}\vec{V}, -{1 \over t}\Hat{V},\\ 
   \sqrt{t}\Hat{V}, -t \vec{V} $$
   この４種類の式から、reco level theory
   が、アダム・スミスの神の見えざる手を経済理論のミクロとマクロ経済の
   具現化を実現させている。
   物理では、$$ {d \over df}F = m(x) = e^f + e^{-f} = 2 i \sin (i x \log x) $$
   と、空間の仕組みが経済理論の場の理論にもなっている。周期関数と
   ピンポイントリサーチが合わさっている。
   これらの式から、前半と後半の株式市場の流れが、シンメトリーとなって
   いるために、後半の$ e^{-f} $が破綻すると、シンメトリーのために、
   前半の$ e^f $も破綻している状態になっている。全体が機能しないと、
   金融の流れが乱れていく。この破綻を回避するには、破綻が連鎖されている
   金融の周期をこれらの式から、調節が可能であることも、類推できる。
   破綻した周期を、正常であった周期に重ね合わせの原理で、ピンポイントリサーチ
   で待って、金融の投資を正常な周期関数で合わせて、元の正常な周期に
   これらの式から戻せることが可能と推測できる。

\title{Abandust judgement on economy strategiest \\
built with acceptant of movement \\
Money demand on law of judgement}
\author{Masaaki Yamaguchi, and investigate from AYA TAKASHIMA \\
also future from my son being have with \\
quantum level on differential geometries}


高島彩さんの研究論文をアカシックレコードと瞑想法で見ての書き出しとしての
論文


   Retry for future to one result with adapt from adulsent for
aquirance of trainning. For must friends of money are treat on
law of judgement, economy need to treat for aquirered of demands,
this demands supply from rest of  one's suppliement.
This balance end without a closed accidennt of out come for
restorance with supplied of demands exceed.
These system are constructed with Actor, Demand, Supply of
being retrayed from Jones manifold around economy of zone
with accesority, verisity of curve in up, down, eternal
rout of stages.

    This system also monument with partial and economy of
similist from law of judgement in universe and money, economy,
human society, estrand of moment in last eternal stage.
  

   Jones manifold equation built with Euler equation and Euler product,
this concernd with concept from Higgs and partial fields,
moreover this concernd of circle and Euler product from
Global integrate and Volume manifold in Sheap manifold of integrate
with group and toplogy different equation.

$$ e^{- \theta} + e^{i \theta} = \sin (i x \log x) = e^{-f} + e^{f} $$
$$ + \cos (i x \log x) = e^{f} - e^{-f} $$

  自動車における道路での速度とそれに対しての加速度を
  平地と上り坂、下り坂における曲がり角での曲率$R_{ij}$が、
  広中平祐先生の庭園理論と同じテーゼで、彩さんが
  経済理論の景気回復と景気刺激剤、景気恐慌が
  どのようにして、自動車運転と自動車の交通規則と
  同じ理論で作用するかを述べているのが、この論文で説明されている。 


  $$ {d \over dx}\left({d \over d r}R_{ij} = \sin
	  \begin{vmatrix} {l_2 \over {2 \pi}} 
- {l_1 \over {2 \pi}}\end{vmatrix}  < 1 \right) = 0 $$


      Curveの度合い $R_{ij}$、潮汐力の差、${d \over d f}F$とは、
   直接的には違うが、間接的には同じであり、rとRの大域的変数
   として、半径rの多様体が曲率的には$R_{ij}$に大域的微分で作用をしている。
   遠隔的に作用している。


      グリーシャ先生からリッチ・フロー方程式が、大域的微分になるのを
   教わった。広中平祐先生の4重帰納法で、ガンマ関数における大域的
   微分多様体が、オイラーの定数の多様体積分に使えるのがわかった。
   これがヒルベルト空間としてのサーストン多様体が、
   代数多様体としての微分幾何として、ゼータ関数が素数分布として
   展開されるのがわかった。

  $$ R_{ij} = |R_2 - R_1| < 1, |R_2 - R_1| = 0 $$
  $$ R_{ij} = {R_2 \over R_1} < 1, {R_2 \over R_1} = 1 $$
  $$ {d \over dt}g_{ij} = - 2 R_{ij}, {d \over d f} F(x) = 
  \int \Gamma(\gamma)^{'}dx_m = \int \Gamma dx_m + {d \over d \gamma}
  \Gamma \le e^{f} - e^{-f} \le e^{-f} + e^{f} $$

\includegraphics[width=6cm,clip]{curve11}
\includegraphics[width=6cm,clip]{curve112}
\includegraphics[width=6cm,clip]{curve12}
\includegraphics[width=6cm,clip]{curve121}
\includegraphics[width=6cm,clip]{curve1211}
\includegraphics[width=6cm,clip]{curve122}

$ R_2 < R_1 \text{の場合は、}$
左からR1に接近して、R1へ行くが、${R_2 \over R_1} < 1 $
より、曲率$R_{ij} < 1$であるから、平坦の道路では、気をつけるべし、
曲率$R_{ij} = 1$では、CURVEには気をつける。
曲率$R_{ij} > 1$では、$R_{i > j}$であり、iの方で気をつける。
これを経済の日経平均相場で、景気の上がりでは、curveがどう分類されるかで、
気をつけるべき時期がわかり、景気の下がりでは、同じくcurveでどう分類するかを
きめる。平坦のときが一番気をつける、油断できない景気の税と立法での
判断が、自動車の走行分岐での心理作用が、景気判断のどの道路状況で
その場の気をつける心理と同じ作用をする。
経済の需要が供給を下回っている状況では、$\text{demand} \le \text{supply}$
では、需要のために、物価を下げるが、この状況では、$ R_1 < R_2, $
であるので、株は$R_2$の方へ曲率がまがり、このときに、需要の
企業の生産が減ってきて、このときに、金融機関が気をつけないと、
企業が破綻する可能性がある。このシチュエーションが世界恐慌であった。
このように、曲率の場合分けをしていないと、間違った判断で、
曲率に気をつけないから、交通事故が起こるということである。
広中平祐先生の庭園理論と同じ考えのモチーフを高島彩さんは、
研究論文で提出している。

   高島彩さんの成蹊大学での研究論文を参照させてもらえて、
書き出した経済の需要と供給に対しての物価上昇と経済沈滞
においてのインフレとデフレにメスを入れる税と景気対策
のタイミング、このタイミングがJones多項式の周期の終わり
にバランスを崩すと、景気の極率としてのCURVEの潮汐力での
差分率の曲率が１以下になるときの、心理作用を考えるべき
という、経済消費における物価調整の作用をかけるメカニズム
についての論文 　

参考文献: 父、母、彩さん、ナッシュ先生、益川先生、ワインバーグ先生、まどかさん


\end{document}
